% TOP_PROBLEM: {"id": 30007267, "title": "The joint Dickman law for consecutive integers", "category_id": 1, "category": {"id": 1, "name": "number_theory", "display_name": "Number Theory", "description": "Properties of integers, prime numbers, Diophantine equations.", "slug": "number-theory", "order_index": 1, "created_at": "2026-07-31T15:26:25.670Z"}, "status": "open"}
% FIRST_POSTED: 2026-10-06
% ATTEMPT_STATUS: solved
% OPENAI_CLAIM: {"schema_version":1,"source_repo":"https://github.com/openai/math","source_commit":"adc7f1241b42e322a6451854ab7e4b4c146bf78a","release_date":"2026-10-06","match":"full","resolution":"proved","independently_reviewed":false,"second_pass":"agreed","families":[{"family":"012","title":"Independent largest prime factors of consecutive integers","subject":"Number theory","manuscripts":[{"title":"The joint Dickman law for consecutive integers","pdf_path":"preprints/The-joint-Dickman-law-for-consecutive-integers-September-24-2026/paper.pdf","readme_path":"preprints/The-joint-Dickman-law-for-consecutive-integers-September-24-2026/README.md","date":"2026-09-24","theorem_ref":"preprints/The-joint-Dickman-law-for-consecutive-integers-September-24-2026/build/sections/introduction.tex :: thm:main"}],"lean":{"doc_path":"lean/docs/012.md","comparators":["lean/ComparatorChallenges/JointDickman.json","lean/ComparatorChallenges/JointDickman.lean"],"declarations":[],"covers_main_theorem":true},"reasoning_trace":null}]}
\documentclass[11pt]{article}
\usepackage[margin=25mm]{geometry}
\usepackage{fontspec}
\setmainfont{Latin Modern Roman}
\usepackage{amsmath,amssymb,mathtools}
\usepackage{unicode-math}
\setmathfont{Latin Modern Math}
\usepackage{enumitem}
\usepackage{xurl,hyperref,bookmark}
\hypersetup{colorlinks=true,urlcolor=blue}
\setcounter{secnumdepth}{0}
\setlength{\parindent}{0pt}
\setlength{\parskip}{5pt}
\setlength{\emergencystretch}{3em}
\begin{document}
\section*{The joint Dickman law for consecutive integers}

\subsection{OpenAI's claim}
OpenAI presents the following claims relevant to this question, with manuscript numbers following paper order in Sources. In manuscript 1 (family 012; 2026-09-24), the asserted result is that consecutive integers have independent limiting largest-prime-factor statistics with Dickman marginals, and either ordering of their largest factors has density one half.

\subsection{Outline of the argument}
For manuscript 1 (family 012; 2026-09-24), a positive divisor weight amplifies a putative mixed correlation into quadratic energy. An independent split model controls its first two moments, while residue smoothing suppresses nonconstant modes and forces the amplified energy to vanish.

\subsection{Scope relative to this problem}
The result supplies the conclusion in this problem with its displayed hypotheses and quantifiers. The broader statements and applications in the manuscript are kept separate from the particular question above. The supporting result is manuscript 1 (family 012; 2026-09-24).

\subsection{Formal verification}
For manuscript 1 (family 012; 2026-09-24), OpenAI lists Lean statements for this assertion: consecutive integers have independent limiting largest-prime-factor statistics with Dickman marginals, and either ordering of their largest factors has density one half.

The central assertion is within the documented formal scope; the additional applications retain their own scope limits.

The linked family Comparator checks are \texttt{JointDickman.lean}.

Their named declarations include \texttt{OAI.JointDickmanPaper.decreasing\_order}, \texttt{OAI.JointDickmanPaper.increasing\_order}, \texttt{OAI.JointDickmanPaper.joint\_law}.

This site has not executed the Lean code. Described statement coverage is not an independent proof check.

\subsection{Status caveat}
The manuscripts were produced by an internal OpenAI model. OpenAI cautions that results without formalization can contain mistakes. This site has not independently checked the proofs; the label records the scope of OpenAI's claim.

\subsection{Sources}
\begin{itemize}
\item OpenAI, \emph{The joint Dickman law for consecutive integers}, 2026-09-24; citation key \texttt{OAI:The-joint-Dickman-law-for-consecutive-integers-September-24-2026}. \href{https://github.com/openai/math/blob/adc7f1241b42e322a6451854ab7e4b4c146bf78a/preprints/The-joint-Dickman-law-for-consecutive-integers-September-24-2026/paper.pdf}{PDF (GitHub)}; \href{https://github.com/openai/math/raw/adc7f1241b42e322a6451854ab7e4b4c146bf78a/preprints/The-joint-Dickman-law-for-consecutive-integers-September-24-2026/paper.pdf}{PDF (direct)}; \href{https://github.com/openai/math/blob/adc7f1241b42e322a6451854ab7e4b4c146bf78a/preprints/The-joint-Dickman-law-for-consecutive-integers-September-24-2026/README.md}{Manuscript page and citation}; \href{https://github.com/openai/math/blob/main/preprints/The-joint-Dickman-law-for-consecutive-integers-September-24-2026/paper.pdf}{Latest version}.
\item Family 012: \href{https://github.com/openai/math/blob/adc7f1241b42e322a6451854ab7e4b4c146bf78a/lean/docs/012.md}{Lean scope}.
\item \href{https://github.com/openai/math/blob/adc7f1241b42e322a6451854ab7e4b4c146bf78a/lean/ComparatorChallenges/JointDickman.json}{Comparator JointDickman.json}.
\item \href{https://github.com/openai/math/blob/adc7f1241b42e322a6451854ab7e4b4c146bf78a/lean/ComparatorChallenges/JointDickman.lean}{Comparator JointDickman.lean}.
\item \href{https://github.com/openai/math/blob/adc7f1241b42e322a6451854ab7e4b4c146bf78a/CONTENTS.md}{OpenAI catalogue entry}.
\end{itemize}
\end{document}
